\documentclass[11pt]{article}

\usepackage[margin=1in]{geometry}
\usepackage{amsmath,amssymb}
\usepackage{booktabs}
\usepackage{graphicx}
\usepackage{xcolor}
\usepackage[colorlinks=true,linkcolor=blue,citecolor=blue,urlcolor=blue]{hyperref}
\usepackage{enumitem}

% Resolve the time-to-target figure from the report directory (relative to docs/).
\graphicspath{{../../outputs/agentic_gs_phase1_reports/hotdog_time_to_target_v7_best/}{./}}

\newcommand{\Nref}{N_{\mathrm{ref}}}
\newcommand{\Ndense}{N}

\title{Agentic3DGS:\\
Reinforcement-Learned Control of 3DGS Training for Acceleration}
\author{Phase~1 technical description}
\date{\today}

\begin{document}
\maketitle

\begin{abstract}
We cast the optimization schedule of 3D Gaussian Splatting (3DGS) as a sequential
decision problem and train a reinforcement-learning (RL) policy to control it
block by block. The policy observes the state of an in-progress reconstruction
(loss statistics, held-out validation quality, compute cost, and the distribution
of Gaussian primitives) and chooses, once per block of optimizer iterations, how
to densify, prune, reset opacity, schedule learning rates, and when to stop. The
reward is a time-discounted, model-complexity-aware measure of held-out quality
improvement, designed so that maximizing return is equivalent to maximizing early
progress on the quality-versus-wall-clock curve---that is, to \emph{accelerating}
training. On the held-out \texttt{hotdog} scene, the learned policy reaches every
test-PSNR target in the $25$--$33$\,dB range in $1.24$--$1.37\times$ less training
wall-clock time than the fixed 3DGS schedule, while keeping the model at or below
the baseline Gaussian budget.
\end{abstract}

\section{Problem Setting}

Stock 3DGS trains a static scene by gradient descent on a photometric loss while
periodically \emph{densifying} (splitting and cloning Gaussians in
high-gradient regions), \emph{pruning} (removing low-opacity or oversized
primitives), and \emph{resetting} opacities. The densification cadence, the
gradient and opacity thresholds, the learning-rate schedule, and the stopping
iteration are all fixed hyperparameters of a hand-tuned schedule. This schedule is
not adaptive: it applies the same recipe regardless of how a particular scene is
converging or how the primitive population is evolving.

We replace the fixed schedule with a learned controller. The 3DGS optimizer is
left untouched and is wrapped in an environment that exposes its internal state
and accepts control actions once per \emph{block} of iterations. A block is a
contiguous run of $B$ optimizer steps; the policy acts at block boundaries, so a
single episode of $T$ iterations consists of roughly $T/B$ decisions rather than
$T$. This block-wise abstraction keeps the action rate low enough for stable RL
while still allowing the controller to adapt within a single training run.

\section{Markov Decision Process}

We model one training run on one scene as an episode of a Markov decision process
$(\mathcal{S}, \mathcal{A}, P, r, \gamma)$. At block $k$ the policy
$\pi_\theta(a_k \mid s_k)$ observes state $s_k$, emits action $a_k$, the
environment executes $B_k$ optimizer iterations under that action, and returns the
next state $s_{k+1}$ and a scalar reward $r_k$. The episode terminates when the
iteration budget is exhausted, when the policy chooses to stop, or when a safety
limit is hit. All quality signals used for state and reward are computed on a
\textbf{held-out subset of training cameras}; the dataset test cameras are never
read during state construction or reward computation. They are used only for
offline reporting and for checkpoint selection (Section~\ref{sec:selection}).

\subsection{State space}
\label{sec:state}

The observation is a $45$-dimensional vector, every component clipped to $[0,1]$.
It groups into:

\begin{description}[leftmargin=1.5em,style=nextline]
  \item[Progress (4)] normalized iteration, elapsed wall-clock fraction,
  remaining-budget fraction, block index.
  \item[Loss (4)] exponential moving averages of the $L_1$ and D-SSIM terms, the
  total block loss, and the recent loss slope (squashed through $\tanh$).
  \item[Validation (4)] held-out PSNR (normalized), SSIM, an LPIPS slot, and the
  per-block quality improvement.
  \item[Compute (5)] time per iteration, time per validation view, peak and
  allocated VRAM fractions, and the active resolution scale.
  \item[Gaussian population, scalar (3)] count fraction, growth rate, and budget
  utilization.
  \item[Previous action (14)] the last block's decoded actions plus a flag for
  whether they improved validation quality, giving the policy a recurrent handle
  on its own behavior.
  \item[Gaussian population, distributional (11)] opacity quantiles
  $q_{10}, q_{50}, q_{90}$, near-transparent fraction, scale quantiles, median
  anisotropy, visible fraction, added/pruned fractions in the previous block, and
  the active spherical-harmonic degree fraction.
\end{description}

The distributional statistics are what let the controller reason about
\emph{health} of the reconstruction---e.g.\ a large near-transparent fraction
signals prunable clutter, high anisotropy signals degenerate splats---rather than
only about scalar loss.

\subsection{Action space}
\label{sec:action}

Each decision is a hybrid action with six discrete heads ($19$ categories total)
and seven continuous components:

\begin{center}
\begin{tabular}{@{}lll@{}}
\toprule
\textbf{Discrete head} & \textbf{Choices} & \textbf{Effect} \\
\midrule
\texttt{block\_steps}            & $\{50,100,200\}$                                  & length of the next block \\
\texttt{densify\_mode}           & off / conservative / default / aggressive          & scales the gradient threshold \\
\texttt{densification\_interval} & $\{50,100,200,400\}$                              & iterations between densifications \\
\texttt{prune\_mode}             & off / opacity / opacity+size                       & pruning criterion \\
\texttt{opacity\_reset}          & none / if-plateau / force                          & opacity-reset trigger \\
\texttt{stop}                    & continue / stop                                    & early termination \\
\bottomrule
\end{tabular}
\end{center}

\begin{center}
\begin{tabular}{@{}lll@{}}
\toprule
\textbf{Continuous component} & \textbf{Range} & \textbf{Scaling} \\
\midrule
\texttt{densify\_threshold\_mult}   & $[0.25, 4.0]$   & log \\
\texttt{prune\_opacity\_threshold}  & $[0.001, 0.02]$ & linear \\
\texttt{position\_lr\_mult}         & $[0.25, 4.0]$   & log \\
\texttt{feature\_lr\_mult}          & $[0.25, 2.0]$   & log \\
\texttt{opacity\_lr\_mult}          & $[0.25, 2.0]$   & log \\
\texttt{scaling\_lr\_mult}          & $[0.25, 2.0]$   & log \\
\texttt{rotation\_lr\_mult}         & $[0.25, 2.0]$   & log \\
\bottomrule
\end{tabular}
\end{center}

Continuous outputs are squashed through a sigmoid and mapped onto their range
(log- or linear-spaced); a neutral action reproduces the stock 3DGS
hyperparameters, so the policy learns \emph{deviations} from a known-good default.

\section{Reward}
\label{sec:reward}

Let $Q_k$ be the held-out validation quality after block $k$, defined per view as
$\mathrm{PSNR} + w_{\mathrm{ssim}}\,\mathrm{SSIM}$ and aggregated across the
held-out cameras as
\begin{equation}
Q_k = \bar{q}_k + w_{\min}\,(q^{\min}_k - \bar{q}_k),
\end{equation}
a blend of the mean and the worst view ($w_{\min}=0.25$). The worst-view term
makes the proxy robust to the policy ``trading away'' poorly covered directions
that the full-hemisphere test cameras would later expose. Let $\Delta Q_k = Q_k -
Q_{k-1}$ be the per-block quality gain.

\paragraph{Model-complexity budget.}
Let $\Ndense$ be the current Gaussian count and $\Nref$ a per-scene reference
budget (set to the initial point-cloud size, $\approx\!10^5$). The over-budget
excess and a smooth complexity discount are
\begin{equation}
e_k = \max\!\Big(0,\ \tfrac{\Ndense}{\Nref}-1\Big), \qquad
d_k = \frac{1}{1 + \kappa\,e_k}, \qquad \kappa = 2.
\end{equation}

\paragraph{Time value of progress.}
The central term for acceleration weights quality gains by the cumulative
\emph{training} wall-clock time $t_k$ already spent:
\begin{equation}
\nu_k = \exp\!\big(-t_k / \tau\big), \qquad \tau = 45\,\text{s}.
\end{equation}
A plain quality reward telescopes to $Q_{\text{final}} - Q_{\text{initial}}$ and
is therefore indifferent to \emph{when} quality is earned; it cannot prefer a
faster schedule. Multiplying each gain by $\nu_k$ makes early progress worth more,
so maximizing return maximizes the early area under the quality-versus-time curve,
which is exactly what a time-to-target metric measures. Early stopping emerges
endogenously: once the discounted marginal gain falls below the marginal time
cost, continuing is no longer rewarded.

\paragraph{Budget-gated growth and compaction.}
Below budget, growth and pruning are both reward-neutral---the objective is
\emph{complexity $\le$ baseline}, not minimal count. Only the over-budget portion
of any change is priced:
\begin{align}
g^{+}_k &= \max\!\big(0,\ \Ndense_k - \max(\Ndense_{k-1}, \Nref)\big)
          &&\text{(growth above budget)}, \\
g^{-}_k &= \max\!\big(0,\ \Ndense_{k-1} - \max(\Ndense_k, \Nref)\big)
          &&\text{(reduction above budget)}.
\end{align}

\paragraph{Total reward.}
The block reward combines a discounted quality credit (positive gains discounted
by complexity, drops kept at full magnitude), a compute-time penalty, explicit
over-budget penalties, a compaction bonus, and stability/resource penalties:
\begin{align}
r_k = \;& \underbrace{w_q\,\nu_k\big(d_k\,[\Delta Q_k]^{+} + [\Delta Q_k]^{-}\big)}_{\text{time-discounted quality}}
        \;-\; w_t\,\frac{\Delta t_k}{t_{\mathrm{ref}}}
        \;-\; w_c\,e_k \nonumber\\
      &\;-\; w_g\,\frac{g^{+}_k}{\Nref}
        \;+\; w_b\,\frac{g^{-}_k}{\Nref}\,\mathbb{1}[\Delta Q_k \ge -\epsilon]
        \;-\; w_v\,[\,\mathrm{VRAM}_k - V^{\star}]^{+}
        \;-\; w_i\,[-\Delta Q_k]^{+},
\end{align}
where $[\,\cdot\,]^{+}$ and $[\,\cdot\,]^{-}$ denote positive and negative parts.
At the terminal block an extra bonus
$w_{tq}\,\nu_k\,d_k\,(Q_k - Q_0) - w_{tc}\,e_k$ rewards total quality gain net of
final excess, and a numerical-failure penalty guards against divergence. Rewards
are standardized per scene (running z-score) before being fed to the RL update, so
that scenes with different natural PSNR scales contribute comparably.

\section{Safety Rails}
\label{sec:safety}

The action range is strictly larger than stock 3DGS, which creates failure modes a
fixed schedule never encounters. The dominant one is a \emph{prune cliff}: stock
3DGS resets opacities to $0.01$ and prunes below $0.005$, so survivors recover, but
the policy may select a prune threshold above the reset value, wiping the model to
the safety floor in a single block. Because validation quality can momentarily
\emph{rise} during such a wipe (it removes near-transparent clutter), the reward
does not see the damage, and a near-random early policy reliably trips the cliff
before learning anything. Two environment-level guards close it---these are
constraints on dynamics, not incentives, and reward tuning alone cannot replace
them:
\begin{enumerate}[leftmargin=1.5em]
  \item \textbf{Post-reset prune guard.} For a recovery window of $300$ iterations
  after any opacity reset, the effective prune threshold is capped at the stock
  value $0.005$.
  \item \textbf{Per-event prune cap.} A single densify/prune event may remove at
  most $25\%$ of the current population, so reaching the floor requires sustained
  pruning across blocks that the reward can observe and penalize.
\end{enumerate}
Additional rails bound opacity-reset frequency, an absolute Gaussian floor, peak
VRAM, and a minimum iteration count before the \texttt{stop} action is honored.

\section{Policy and Optimization}

\paragraph{Architecture.}
A shared multilayer perceptron (3 hidden layers of width $256$, GELU activations)
encodes the state. From the shared features, six independent categorical heads
produce the discrete-action logits, a diagonal Gaussian head (learned log-std)
produces the seven continuous actions, and a scalar critic head produces the
value estimate.

\paragraph{Training.}
The policy is optimized with Proximal Policy Optimization (PPO): clipped surrogate
objective (clip $0.2$), generalized advantage estimation ($\gamma=0.99$,
$\lambda=0.95$), advantage normalization, value clipping, gradient-norm clipping,
an entropy bonus, and KL-based early stopping of the inner epochs. Rollouts span
multiple episodes across a diverse set of training scenes
(\texttt{chair}, \texttt{drums}, \texttt{mic}, \texttt{ficus}, \texttt{ship});
scene diversity is essential, since a controller trained only on easy scenes
overfits to their capacity profile and transfers poorly. The hybrid log-probability
is the sum of the categorical log-probabilities and the Gaussian
log-probabilities, and the entropy bonus sums both, so exploration is driven
jointly over discrete and continuous decisions.

\section{Acceleration-Aware Checkpoint Selection}
\label{sec:selection}

Because the train-camera reward is a proxy, over-optimizing it can widen the
train-/test-camera gap; selecting checkpoints by reward or by update count
reproduced a regression in which \emph{more} training yielded \emph{worse} test
PSNR. We therefore select checkpoints with a held-out \emph{probe}: every few PPO
updates, the current policy runs one frozen, deterministic episode on a held-out
probe scene (neither a training nor an evaluation scene), scored on that scene's
\textbf{test} cameras. Crucially, the probe is scored for \emph{acceleration}: test
PSNR is sampled on a fast view subset whenever cumulative training wall-clock time
crosses $\{15, 30, 60\}$\,s, and the selection metric is the mean of these
time-sliced PSNRs---an approximation to the early area under the quality-time
curve. The checkpoint maximizing this score is saved as the deployed policy. With
this metric the objective improves monotonically over training, the opposite of
the reward-/count-selected regression.

\section{Results}

We evaluate with a time-to-target protocol: both the fixed 3DGS schedule and the
frozen policy are run under matched iteration budgets and the same camera split,
test PSNR is measured on the full test set, and the training wall-clock time to
reach each target PSNR is read off the PSNR-time Pareto envelope (linearly
interpolated between samples). On the held-out \texttt{hotdog} scene:

\begin{center}
\begin{tabular}{@{}rrrr@{}}
\toprule
\textbf{Target PSNR (dB)} & \textbf{Baseline (s)} & \textbf{Agentic (s)} & \textbf{Speedup} \\
\midrule
25 & 9.72  & 7.08  & $1.37\times$ \\
27 & 14.17 & 11.34 & $1.25\times$ \\
29 & 18.27 & 13.91 & $1.31\times$ \\
31 & 22.37 & 17.99 & $1.24\times$ \\
33 & 30.98 & 24.22 & $1.28\times$ \\
\bottomrule
\end{tabular}
\end{center}

The learned controller dominates the fixed schedule across the target range at a
constant, baseline-sized Gaussian budget, reaching each quality level in
$1.24$--$1.37\times$ less training time. The policy stops itself at roughly $3000$
iterations (a self-discovered consequence of the time-discounted reward), which
caps its frontier near $34$\,dB; pushing the high-PSNR range further is a matter of
lengthening the time-value horizon $\tau$. Figure~\ref{fig:ttp} shows the full
quality-versus-time curves.

\begin{figure}[t]
  \centering
  % Generated by scripts/plot_time_to_target_artifacts.py into the report dir:
  % outputs/agentic_gs_phase1_reports/hotdog_time_to_target_v7_best/time_to_psnr.pdf
  \includegraphics[width=0.82\textwidth]{time_to_psnr.pdf}
  \caption{Held-out test PSNR versus training wall-clock time on \texttt{hotdog}.
  The agentic policy (orange) reaches every quality level to the left of the fixed
  3DGS schedule (blue), then stops itself near $34$\,dB; the baseline continues to
  a higher ceiling at greater cost.}
  \label{fig:ttp}
\end{figure}

\section{Summary}

The contribution is a control reformulation of 3DGS training in which an RL policy
adaptively schedules densification, pruning, opacity resets, learning rates, and
stopping, driven by a reward whose maximizer is fast early convergence at a bounded
model size. Three design elements make it work: budget-gated complexity terms that
make ``complexity $\le$ baseline'' the operative constraint without forcing minimal
count; dynamics-level safety rails that remove an otherwise unavoidable model-wipe
failure mode; and an acceleration-aware probe that selects checkpoints by
time-to-quality on held-out test cameras rather than by training reward. Together
they yield a controller that accelerates held-out reconstruction by
$1.24$--$1.37\times$ at matched quality and budget.

\end{document}
